\section{Clearing a small pending component}\label{sec:clearing}

An ordinary sweep can leave active sites marked pending. Their partial
potentials may contain auxiliary digits, and a failed trial may already have
exposed true digits that cannot be used again as fresh inputs. A successful clearing
completes a small component of pending sites to its true values and removes
its damping. The estimate for this completion must hold for the true values
obtained, without a freshness assumption on the core.

The preparation uses true residuals at selected exterior sites. To explain
the selection, we first identify the directions through which completing the
core can change its exterior. This gives a projection bound that suffices for
completion. We then show that the whole relevant range is spatially localized,
and construct the projection bound by a stopped sequence of exterior exposures.
The result in this section is conditional on explicit local geometry and
residual laws. Section~\ref{sec:assembly} verifies those inputs for the
scheduled pending components and combines the successful changes.

Fix the dimension $d\ge3$. Throughout this section the damping scale $t$
and the vector $r$ are fixed; constants may depend on $d$, but not on the
torus volume or on the particular component. The current active inverse is
\[
 g=(M-t\diag r)^{-1},\qquad M=M^\top\in\RR^{A\times A},
 \qquad h=\sqrt t\,g,
\]
where $M$ includes the Schur complement of the true terminal sites. Write
$p_i=\Im r_i$ and assume the fixed broad bounds
\begin{equation}\label{eq:clearing-broad}
 |r_i|\le C,\qquad c\le p_i\le C.
\end{equation}
Split the active sites into a fixed core $B$ and its current exterior
$R=A\setminus B$. During preparation, exterior sites can be removed from
the trial's active set. The direct partial potential values on $B$ stay
uncompleted until the final core change, but eliminating an exterior site
changes inverse entries involving $B$ and hence the current scattering
cross block.

We use the parameters
\begin{equation}\label{eq:clearing-parameters}
\begin{gathered}
 \delta=.01,\qquad \nu=10^{-5}\delta/d,\\
 \beta_w=.36\delta,\qquad e_0=.44\delta,\qquad e_h=.46\delta,\\
 H_*=t^{1+\delta},\qquad r_*=t^{-\nu},\qquad
 l_*=t^{1+e_0},\qquad n_h=\lfloor t^{-1-e_h}\rfloor.
\end{gathered}
\end{equation}
The empty core needs no change. For the main development assume
$1\le |B|\le r_*$; the stopped trial below includes the empty case.

For a terminal set $T$, give every nearest-neighbor edge touching $T$ cost
zero and every other such edge cost one. Let $d_T$ be the resulting
shortest-path pseudometric. A nonlocal entry created by eliminating a
connected terminal component joins active sites at zero $d_T$-distance.
For a set $B$, put $d_T(i,B)=\min_{b\in B}d_T(i,b)$.

\subsection{The scattering range and the completion target}

Write $D_p=\diag(p)$ and define
\begin{equation}\label{eq:clearing-scattering}
 S=\Id+2itD_p^{1/2}gD_p^{1/2}.
\end{equation}
The matrix $S$ is symmetric and unitary. Indeed, for
$Z=D_p^{1/2}gD_p^{1/2}$ the Ward identity gives
$\operatorname{Im}Z=tZZ^*$, and hence
\[
 SS^*=\Id+2it(Z-Z^*)+4t^2ZZ^*=\Id.
\]
Write
\[
 \mathcal C=S_{RB},\qquad \Pi=\Pi_{\Ran\mathcal C}
\]
for the scattering cross block and the orthogonal projection onto its
range in $\CC^R$. Its diagonal has the concrete meaning
\[
 \Pi_{ii}=\sup_{v\in\Ran\mathcal C,\ \|v\|=1}|v_i|^2,
\]
with value zero when the range is zero. Thus a small diagonal means that
no unit vector in the cross-block range can concentrate at one exterior
site. The next lemma shows why this is the condition needed to complete
the core.

\begin{lemma}[Core completion from its scattering range]\label{lem:core-completion}
Consider an active split $B,R$, with $R\ne\varnothing$, frozen positive
bounded $p$, damping $t$, and $|B|\le r_*$. Suppose every relevant true
terminal principal matrix is invertible and the projection $\Pi$ onto
$\Ran\mathcal C$, where $\mathcal C=S_{RB}$, satisfies $\Pi_{ii}\le2l_*$.
Set the potential on $B$ to its true value, remove its damping, and
eliminate $B$. The new survivor inverse $g'$ exists. Write
$\Delta g=g'-g_{RR}$, $\Delta h=\sqrt t\,\Delta g$, and
$\Delta S=S'-S_{RR}$, where $S'$ is the scattering matrix of the new
survivor inverse. Then
\[
 \|\Delta S\|\le2,\qquad
 \Delta S=\Pi\Delta S\overline\Pi,
\]
and
\begin{align}
 \max_{i,j\in R}|\Delta g_{ij}|&\le Ct^{e_0},
 &\max_{i\in R}\sum_{j\in R}|\Delta h_{ij}|^2&\le Ct^{e_0},
 \label{eq:assembly-core-rows}\\
 \max_{i\in R}\sum_{j\in R}|\Delta h_{ij}|^4&\le Ct^{1+3e_0},
 &\sum_{i,j\in R}|\Delta h_{ij}|^2&\le Cr_*/t.
 \label{eq:assembly-core-fourth}
\end{align}
These bounds are deterministic for every true core completion satisfying
the stated invertibility hypothesis.
\end{lemma}

\begin{proof}
Let $T_0$ denote the existing terminal set. The true principal matrices on
$T_0$ and on $T_0\cup B$ are invertible.  Eliminating the latter set therefore
produces a real symmetric survivor matrix with damping $t\diag(p_R)$,
whose inverse exists by dissipativity.  Reversing these Schur eliminations
shows that the full active matrix, with the new diagonal on $B$, is
invertible as well.

Let $V$ be its diagonal change on $B$; its entries include the removal
of $-tr_i$.  In particular $\Im V_{ii}=tp_i>0$ and $V$ is invertible.
The finite-dimensional resolvent identity gives
\begin{equation}\label{eq:clearing-core-resolvent}
 \Delta g=-g_{RB}(\Id+Vg_{BB})^{-1}Vg_{BR}.
\end{equation}
To see that the middle inverse exists, write $\mathcal M=g^{-1}$ for
the old active matrix and $E_B:\CC^B\to\CC^{B\cup R}$ for coordinate
inclusion.  The finite-rank determinant identity gives
\[
 \det(\mathcal M+E_BVE_B^\top)
       =\det(\mathcal M)\det(\Id+Vg_{BB}),
\]
and both full determinants are nonzero.  Thus no estimate on the norm of
this middle inverse is assumed.

The old scattering matrix $S$ is symmetric unitary.  Its $R$ block
$S_{RR}$ is a contraction, whereas the new survivor scattering matrix $S'$
is symmetric unitary.  Hence $\Delta S=S'-S_{RR}$ satisfies
$\|\Delta S\|\le2$.
Since $g_{BR}=g_{RB}^\top$, \eqref{eq:clearing-core-resolvent} yields
\[
 \Delta S=\mathcal C C_*\mathcal C^\top,\qquad \mathcal C=S_{RB},
\]
where the coefficient $C_*$ incorporates the invertible diagonal $p_B$
factors.  Its size is irrelevant.  Left range inclusion gives
$\Pi\Delta S=\Delta S$.  Transposing this equality and using
$\Delta S^\top=\Delta S$ gives $\Delta S\Pi^\top=\Delta S$.
An orthogonal projection is Hermitian, so $\Pi^\top=\overline\Pi$; thus
\begin{equation}\label{eq:clearing-core-support}
 \Delta S=\Pi\Delta S\overline\Pi.
\end{equation}
The conjugation on the right is necessary for a generally complex range.
The two projections have the same real diagonal.  Cauchy--Schwarz and the
operator norm bound imply
\[
 |\Delta S_{ij}|\le2\sqrt{\Pi_{ii}\Pi_{jj}}\le4l_*.
\]
Similarly, $\Delta S\Delta S^*\le4\Pi$ gives
$\sum_j|\Delta S_{ij}|^2\le4\Pi_{ii}\le8l_*$.
The factorization has rank at most $|B|$, whence
$\|\Delta S\|_{\mathrm{HS}}^2\le4|B|$.

Finally,
\[
 \Delta S=2i\sqrt t\,\sqrt{p_R}\,\Delta h\,\sqrt{p_R}.
\]
Bounded positive $p$ gives respectively
$\max|\Delta h_{ij}|\le Cl_*/\sqrt t=C\sqrt t\,t^{e_0}$,
$\max_i\sum_j|\Delta h_{ij}|^2\le Cl_*/t=Ct^{e_0}$, and
$\sum_{ij}|\Delta h_{ij}|^2\le C|B|/t$.
The fourth-power row sum is at most the largest squared entry times the
row square sum, namely $Ct^{1+2e_0}t^{e_0}$.
This proves \eqref{eq:assembly-core-rows}--\eqref{eq:assembly-core-fourth}.
\end{proof}

The middle coefficient in the core resolvent factorization can be large.
The proof controls the resulting scattering change by its operator norm and
its range instead. This explains why the raw row norms of $\mathcal C$ are
insufficient: a small singular amplitude can still correspond to a range
direction concentrated at one exterior site. We must control the full range,
including every nonzero singular direction. The exterior equation gives
that range a spatial bound before any diagonal lag or comparison estimate
is imposed.

\subsection{Localization of the full scattering range}

\begin{lemma}[Localization of the full scattering range]
\label{lem:clearing-range-localization}
Let $B,R$ split the active sites of a matrix obtained by eliminating true
terminal sites, and suppose the active damping is $t\diag(p)$ with
$c\le p_i\le C$.  For $\mathcal C=S_{RB}$ and the orthogonal projection
$\Pi$ onto $\Ran \mathcal C$, there is $a>0$, depending only on the fixed bounds and
on $d$, such that
\begin{equation}\label{eq:clearing-projection-decay}
 \Pi_{ii}\le C\exp\{-2at(d_T(i,B)-1)_+\},\qquad i\in R.
\end{equation}
This assertion requires no bound on $g_{ii}$, $K$, or the inverse energy
form, and no lower bound on the nonzero singular values of $\mathcal C$.
\end{lemma}

\begin{proof}
Write $\mathcal M=M-t\diag(r)$, so $g=\mathcal M^{-1}$.  The exterior block is dissipative:
for every $z\in\CC^R$,
\[
 -\Im\langle z,\mathcal M_{RR}z\rangle
       =t\sum_{i\in R}p_i|z_i|^2\ge ct\|z\|^2.
\]
Cauchy--Schwarz implies $\|\mathcal M_{RR}z\|\ge ct\|z\|$, and hence
$\mathcal M_{RR}$ is invertible with $\|\mathcal M_{RR}^{-1}\|\le C/t$.
For any coefficient vector $z\in\CC^B$, the vector $w=g_{RB}z$ obeys
\begin{equation}\label{eq:clearing-exterior-equation}
 \mathcal M_{RR}w=-\mathcal M_{RB}g_{BB}z=:f.
\end{equation}
Every nonzero off-diagonal entry of $M$ is either a direct lattice hop or
comes from the terminal Schur correction.  The latter correction is block
supported on active neighbors of a connected terminal component: the
terminal matrix and its inverse are block diagonal by those components.
Any two such neighbors have collapsed distance zero.  Consequently the
source $f$ is supported on $\{i:d_T(i,B)\le1\}$.

Let $\rho_i=(d_T(i,B)-1)_+$ and let $\mathcal E$ be multiplication by
$e^{at\rho_i}$.  Then $\mathcal E f=f$.  The function $\rho$ is one-Lipschitz
in the collapsed pseudometric.  The diagonal and every terminal Schur
entry commute with $\mathcal E$, including arbitrarily large Schur entries.
Only direct hopping entries can contribute to
\[
 E_a=\mathcal E \mathcal M_{RR}\mathcal E^{-1}-\mathcal M_{RR}.
\]
For each such entry, $|\rho_i-\rho_j|\le1$, so its change is at most
$e^{at}-1\le C at$ for $t\le1$.  There are at most $2d$ direct neighbors
in each row and column.  The absolute row and column sum bounds therefore
give $\|E_a\|\le C_d at$.
Using \eqref{eq:clearing-exterior-equation} now gives the exact identity
\[
 (\mathcal M_{RR}+E_a)\mathcal Ew=f=\mathcal M_{RR}w,
 \qquad
 \mathcal Ew=w-\mathcal M_{RR}^{-1}E_a\mathcal Ew.
\]
Choose $a>0$ so small that $\|\mathcal M_{RR}^{-1}E_a\|\le1/2$.
Absorption proves
\begin{equation}\label{eq:clearing-range-weight}
 \|\mathcal Ew\|\le2\|w\|\qquad(w\in\Ran g_{RB}).
\end{equation}
All matrices are finite, so no domain issue arises from the exponential
weight.  Crucially, the right side is the norm of $w$ itself, not the norm
of a vector used to represent it.

Because
$\mathcal C=2it\sqrt{p_R}\,g_{RB}\sqrt{p_B}$, the matrix $\sqrt{p_B}$ is
invertible, and $\sqrt{p_R}$ commutes with $\mathcal E$, bounded positive
$p$ transfers \eqref{eq:clearing-range-weight} to
$\|\mathcal Ev\|\le C\|v\|$ for all $v\in\Ran \mathcal C$.
Thus $|v_i|\le C e^{-at\rho_i}\|v\|$ on this entire range.
The squared norm of coordinate evaluation on $\Ran \mathcal C$ is $\Pi_{ii}$;
equivalently,
$\Pi_{ii}=\sup_{v\in\Ran \mathcal C,\,\|v\|=1}|v_i|^2$.
This proves \eqref{eq:clearing-projection-decay}.  When $\mathcal C=0$ the projection
is zero and the same conclusion holds.  Rank deficiency and arbitrarily
small positive singular values never entered the estimate.
\end{proof}

\begin{corollary}[Spatial decay of a core completion]\label{cor:core-decay}
Under the hypotheses of Lemma~\ref{lem:core-completion}, let $d_T$ be the
collapsed distance immediately before completion. The same survivor change
satisfies
\begin{equation}\label{eq:clearing-core-decay}
 |\Delta h_{ij}|\le Ct^{-1/2}
 \exp\{-at[(d_T(i,B)-1)_++(d_T(j,B)-1)_+]\}.
\end{equation}
\end{corollary}

\begin{proof}
In the support identity \eqref{eq:clearing-core-support}, use
Lemma~\ref{lem:clearing-range-localization} at both endpoints in place of
the uniform bound $\Pi_{ii}\le2l_*$. The estimate
$|\Delta S_{ij}|\le2\sqrt{\Pi_{ii}\Pi_{jj}}$ and
$\Delta S=2i\sqrt t\,\sqrt{p_R}\,\Delta h\,\sqrt{p_R}$ give the result.
\end{proof}

The same range estimate will locate the exterior exposures used to reach
the projection bound. We now construct that bound while controlling the
accumulated change of the inverse.

\subsection{Exterior heats and regularized projection weights}
\label{subsec:clearing-algebra}

An exterior update, called a heat step, replaces the
partial potential at one exterior site $x$ by its true value and removes that
site from the trial's active set. Such removal remains provisional until
the clearing is committed; rollback restores the pretrial matrix and partial
potential data while retaining the revealed values in the history.
If its unexposed true residual is $\xi$, the change
of its diagonal in the damped matrix is $\xi+tr_x$.  The conditional hypotheses
on the residual are
\begin{equation}\label{eq:clearing-residual}
 \EE(\xi\mid\mathcal F)=0,\qquad
 \EE(\xi^2\mid\mathcal F)=t-\eta,\qquad
 |\xi|\le C\sqrt t,\qquad \eta/t\le Ct.
\end{equation}
Here $\mathcal F$ is the recorded history before the site is inspected.  The
selection of $x$ must be $\mathcal F$-measurable. For the scheduled
components, Section~\ref{sec:assembly} combines the pivot geometry with
Lemma~\ref{lem:conditional-freshness} to verify these hypotheses.

There are two distinct sets of state bounds in the argument.  At the beginning
of a component trial we have fourth-power row slack and a small diagonal lag.
At a provisional heat prefix we use only
\begin{equation}\label{eq:clearing-prestate}
 \max_i|g_{ii}-r_i|\le C t^{\beta_w},\qquad
 \max_i\sum_j|h_{ij}|^4\le H_*.
\end{equation}
Together with \eqref{eq:clearing-broad}, Ward's identity gives
\begin{equation}\label{eq:clearing-row-bound}
 \sum_jp_j|h_{ij}|^2=\Im g_{ii},\qquad
 \max_i\sum_j|h_{ij}|^2\le C.
\end{equation}
The provisional estimates below use no bound on $K$ and no comparison of the
Laplacian $D$ with a reference form. Section~\ref{sec:assembly} transports
those comparisons for the resulting committed operation.

At each heat the core $B$ remains fixed, but its current exterior $R$ and
the cross block $\mathcal C=S_{RB}$ change. Write $S_B=S_{BB}$ and set
\[
 \Gamma=\mathcal C^*\mathcal C.
\]
Unitarity gives $\Gamma=\Id-S_B^*S_B$.
For a ridge parameter $\tau>0$, define
\begin{equation}\label{eq:clearing-ridge}
 G_\tau=\Gamma+\tau\Id,\qquad
 U=\mathcal C G_\tau^{-1/2},\qquad
 m_i=\|u_i\|^2\quad(i\in R),
\end{equation}
where $u_i$ is the $i$-th row of $U$. Equivalently, $m_i$ is the $i$-th
diagonal entry of
\[
 \mathcal C(\mathcal C^*\mathcal C+\tau\Id)^{-1}\mathcal C^*.
\]
These entries are the standard ridge leverage scores
\cite[Section~3.3, Definition~1]{ElAlaouiMahoney2015}; here we use them to
select predictable exterior residual exposures while the ridge decreases
toward the full range projection.

To see exactly what the regularization retains, write a singular-value
decomposition
$\mathcal C=\sum_{\alpha:\sigma_\alpha>0}\sigma_\alpha v_\alpha w_\alpha^*$,
with orthonormal left and right singular vectors. Then
\[
 \mathcal C(\Gamma+\tau\Id)^{-1}\mathcal C^*
 =\sum_{\alpha:\sigma_\alpha>0}
       \frac{\sigma_\alpha^2}{\sigma_\alpha^2+\tau}
       v_\alpha v_\alpha^*.
\]
As $\tau\downarrow0$, every nonzero singular direction receives weight
one, including directions with arbitrarily small singular value. The
matrix therefore increases to $\Pi_{\Ran\mathcal C}$, and each $m_i$ is
at most the corresponding projection diagonal. No zero singular direction
contributes. For fixed $\mathcal C$, each $m_i$ is continuous for
$\tau>0$ and
\[
 \frac{d}{d\tau}m_i(\tau)
 =-\mathcal C_i(\Gamma+\tau\Id)^{-2}\mathcal C_i^*\le0.
\]
Thus decreasing the ridge increases the regularized weights continuously
toward the completion target.

Put
\[
 m_*=\max_{i\in R}m_i,\qquad
 \mathcal A=\sum_{i\in R}m_i,\qquad
 \mathcal P=\sum_{i\in R}m_i^2.
\]
If $R$ is empty, the maximum is zero. The matrices
\begin{equation}\label{eq:clearing-ridge-defect}
 Q=U^*U,\qquad V_\tau=\Id-Q=\tau G_\tau^{-1}
\end{equation}
satisfy $0\le Q,V_\tau\le\Id$, so
$\mathcal A=\Tr Q\le\operatorname{rank}\mathcal C\le |B|\le r_*$.

Whenever the maximum regularized weight is at least $2l_*$ and at most
$4l_*$, the procedure heats a site of maximum weight, chosen from the
recorded history before its residual is inspected. Removing that pivot deletes its contribution
$m_*^2$ from $\mathcal P$; the issue is to show that redistribution to
the surviving rows is smaller. Maximality also gives $m_i\le m_*$ in
the estimates for that redistribution. The next lemma establishes this
fixed-ridge drift under the provisional prestate bounds.

A heat holds the scalar $\tau$ fixed, changes the damped diagonal at its
chosen site $x$ by $\xi+tr_x$, and provisionally eliminates $x$ as a true
terminal. The standing invertibility of true terminal principal matrices
makes the elimination well-defined. The new cross block is recomputed from
the new active inverse at the same scalar ridge. Use primes for quantities
on $R'=R\setminus\{x\}$.

\begin{lemma}[Conditional leverage drift]\label{lem:leverage-drift}
Suppose the active inverse \(g\), with frozen \(r,p\), satisfies
\[
 |g_{ii}|+|r_i|\le C,\qquad
 c\le p_i\le C,\qquad
 |g_{ii}-r_i|\le Ct^{\beta_w},\qquad
 \sum_j|h_{ij}|^4\le H_*
\]
at the prestate of a heat, where \(h=\sqrt t\,g\).  Suppose
\(|B|\le r_*\), \(\tau>0\), and \(x\) is a predictable maximizer of
the ridge leverage, with
\[
 2l_*\le m_x=m_*\le4l_*.
\]
Conditionally on this prestate, assume that \(\xi\) is real and that
\begin{equation}\label{eq:clearing-residual-moments}
 \EE\xi=0,\qquad \EE\xi^2=t-\eta,\qquad
 |\xi|\le C\sqrt t,\qquad 0\le\eta/t\le Ct.
\end{equation}
If
\begin{equation}\label{eq:clearing-drift-exponents}
 0\le\nu<\frac{1+\delta}{2},\qquad
 0<e_0<\delta-\nu,\qquad
 \frac{e_0}{2}<\beta_w<\frac12,
\end{equation}
then, uniformly over these prestates as \(t\downarrow0\),
\[
 |\EE(\mathcal A'-\mathcal A)|=o(m_*),\qquad
 \EE(\mathcal P'-\mathcal P)\le-(1-o(1))m_*^2.
\]
These conclusions include an update at which a poststate cap first fails;
only the displayed prestate assumptions are required.
\end{lemma}

\begin{proof}
All expectations in the proof are conditional on the prestate.  Write
\(m=m_*=m_x\).  The Sherman--Morrison identity gives, on surviving indices,
\begin{equation}\label{eq:clearing-heat-update}
 h'_{ij}=h_{ij}+c_xh_{ix}h_{xj},\qquad
 c_x=-\frac{\xi+tr_x}
 {\sqrt t\,\{1+(\xi+tr_x)g_{xx}\}}.
\end{equation}
The denominator is bounded away from zero because \(g_{xx}\) is bounded
and \(|\xi+tr_x|\le C\sqrt t\).  Expansion to second order gives
\[
 \EE c_x
 =\sqrt t\,(g_{xx}-r_x)-\frac{\eta}{\sqrt t}g_{xx}+O(t).
\]
Here the remainder follows from
\(\EE|\xi+tr_x|^3\le Ct^{3/2}\).
Since \(\beta_w<1/2\), it follows that
\begin{equation}\label{eq:clearing-heat-moments}
 |c_x|\le C,\qquad
 |\EE c_x|\le C\sqrt t\,t^{\beta_w},\qquad
 \EE|c_x|^2=1+O(\sqrt t+\eta/t).
\end{equation}

For \(i\in R'\), let
\[
 v_i=h_{ix}\sqrt{p_i/p_x}.
\]
The Ward identity and the diagonal lag bound show that
\[
 \sum_{\text{active }i}|h_{ix}|^2\frac{p_i}{p_x}
 =\frac{\operatorname{Im}g_{xx}}{p_x}
 =1+O(t^{\beta_w}).
\]
Each omitted index in \(B\cup\{x\}\) contributes at most
\(C\sqrt{H_*}\).  Thus
\begin{equation}\label{eq:clearing-v-bounds}
 \|v\|^2=1+O(t^{\beta_w}+r_*\sqrt{H_*}),\qquad
 \max_i|v_i|\le CH_*^{1/4},\qquad
 \sum_i|v_i|^4\le CH_*.
\end{equation}
The cross block in the old normalization is, exactly,
\begin{equation}\label{eq:clearing-cross-update}
 W=\mathcal C'G_\tau^{-1/2}=U_{R'}+c_xvu_x.
\end{equation}
There is no complex conjugate on \(u_x\) in this formula, because the
resolvent update is symmetric.

\paragraph{The unitary cancellation.}
Define \(b_x=U_{R'}^*v\).  The estimate needed for \(b_x\) is stronger
than the direct bound \(\|U_{R'}\|\|v\|\).  Orthogonality of the \(B\)
columns of \(S\) to its \(x\)-th column gives
\[
 \mathcal C^*S_{Rx}=-S_B^*S_{Bx}.
\]
Since \(S_B\) is symmetric,
\(\overline\Gamma=\Id-S_BS_B^*\), and
\(\Gamma S_B^*=S_B^*\overline\Gamma\).  Functional calculus therefore
gives the intertwining identity
\[
 G_\tau^{-1/2}S_B^*
   =S_B^*(\overline\Gamma+\tau\Id)^{-1/2}.
\]
Also \(S_{Bx}=\mathcal C_x^\top\), so symmetry gives
\[
 \|(\overline\Gamma+\tau\Id)^{-1/2}S_{Bx}\|
   =\|\mathcal C_xG_\tau^{-1/2}\|=\sqrt m.
\]
As \(\|S_B\|\le1\), we obtain \(\|U^*S_{Rx}\|\le\sqrt m\).
For \(i\in R'\),
\(S_{ix}=2i\sqrt t\,p_xv_i\).  Restoring and then subtracting the
\(x\)-th term in \(U^*S_{Rx}\) yields
\begin{equation}\label{eq:clearing-b-cancellation}
 b_x=\frac{U^*S_{Rx}-u_x^*S_{xx}}{2i\sqrt t\,p_x},
 \qquad \|b_x\|\le C\sqrt{m/t}.
\end{equation}

\paragraph{The Gram increment.}
Set
\[
 D_* = W^*W-Q,\qquad
 \varepsilon_t=t^{\beta_w}+\sqrt t+r_*\sqrt{H_*}.
\]
Multiplication of \eqref{eq:clearing-cross-update} gives
\begin{equation}\label{eq:clearing-gram-increment}
 D_*=(|c_x|^2\|v\|^2-1)u_x^*u_x
       +c_xb_xu_x+\overline{c_x}u_x^*b_x^*.
\end{equation}
Let \(\|\cdot\|_1\) denote the trace norm.  Equations
\eqref{eq:clearing-heat-moments}--\eqref{eq:clearing-b-cancellation}
imply
\begin{equation}\label{eq:clearing-gram-bounds}
 \|D_*\|_1\le C\frac m{\sqrt t},\qquad
 \|\EE D_*\|_1\le Cm\varepsilon_t,\qquad
 \Tr D_*^2\le C\frac{m^2}{t}.
\end{equation}
For example, the conditional mean of the two mixed terms is bounded
by \(2|\EE c_x|\|b_x\|\sqrt m\le Cmt^{\beta_w}\).
The final bound follows from
\(\Tr D_*^2\le\|D_*\|_1^2\).
Since \(m\le4t^{1+e_0}\), we have \(\|D_*\|=o(1)\), pathwise.

The new Gram matrix satisfies the exact congruence
\begin{equation}\label{eq:clearing-gram-congruence}
 \Gamma'+\tau\Id
   =G_\tau^{1/2}(\Id+D_*)G_\tau^{1/2}.
\end{equation}
Indeed, \(W^*W=Q+D_*\) and \(\tau G_\tau^{-1}=\Id-Q\).
Inverting this congruence shows that the new leverages are
\begin{equation}\label{eq:clearing-leverage-update}
 m_i'=w_i(\Id+D_*)^{-1}w_i^*.
\end{equation}
No commutation of \(D_*\) and \(G_\tau\) is used here.

\paragraph{Total leverage.}
By cyclicity of the trace and \(W^*W=Q+D_*=\Id+D_*-V_\tau\),
\begin{equation}\label{eq:clearing-total-leverage}
 \mathcal A'-\mathcal A
 =\Tr\{V_\tau(\Id+D_*)^{-1}D_*\}.
\end{equation}
The identity
\[
 (\Id+D_*)^{-1}D_*
   =D_*-D_*^2(\Id+D_*)^{-1}
\]
and \(0\le V_\tau\le\Id\) imply
\[
 |\EE(\mathcal A'-\mathcal A)|
 \le \|\EE D_*\|_1+C\EE\Tr D_*^2
 \le Cm(\varepsilon_t+t^{e_0})=o(m).
\]
The quadratic remainder is controlled by \(\Tr D_*^2\), without a
factor depending on \(|B|\).

\paragraph{Squared leverage.}
It remains to show that the loss of the pivot row dominates the
redistribution of leverage among the survivors.  Put
\[
 d_i=m_i'-m_i,\qquad
 d_i^0=\|w_i\|^2-\|u_i\|^2
      =2\operatorname{Re}(c_xv_iu_xu_i^*)
        +|c_x|^2|v_i|^2m.
\]
Then
\begin{equation}\label{eq:clearing-square-decomposition}
 \mathcal P'-\mathcal P
 =-m^2+2\sum_{i\in R'}m_id_i+\sum_{i\in R'}d_i^2.
\end{equation}
Since \(\|U\|\le1\),
\[
 \sum_i|u_xu_i^*|^2\le m,\qquad
 \sum_i m_i^2\le m\mathcal A\le mr_*.
\]
Cauchy--Schwarz and \eqref{eq:clearing-v-bounds} therefore give
\begin{equation}\label{eq:clearing-raw-mean}
 \left|\sum_i m_i\EE d_i^0\right|
 \le C|\EE c_x|m\sqrt m+Cm\sqrt{r_*mH_*}.
\end{equation}

To handle normalization, expand
\[
 (\Id+D_*)^{-1}-\Id
   =-D_*+D_*^2(\Id+D_*)^{-1}.
\]
The first-order term with the old rows has conditional mean bounded by
\[
 \left|\Tr\left\{\left(\sum_i m_iu_i^*u_i\right)
                    \EE D_*\right\}\right|
 \le m\|\EE D_*\|_1,
\]
because \(\|\sum_i m_iu_i^*u_i\|\le m\).
Replacing the old rows by the updated rows costs at most
\begin{equation}\label{eq:clearing-row-replacement}
 Cm\sqrt m\,\|D_*\|_1.
\end{equation}
For completeness, if \(M_R=\diag((m_i)_{i\in R'})\), then
\[
 W^*M_RW-U_{R'}^*M_RU_{R'}
 =c_xU_{R'}^*M_Rvu_x
  +\overline{c_x}u_x^*v^*M_RU_{R'}
  +|c_x|^2(v^*M_Rv)u_x^*u_x.
\]
Its trace norm is at most \(Cm\sqrt m+Cm^2\le Cm\sqrt m\),
using \(\|U_{R'}\|\le1\), \(\|M_R\|\le m\), and \(\|v\|\le C\).
Taking its trace against \(D_*\) proves
\eqref{eq:clearing-row-replacement}.
Moreover, \(\|W\|\le2\) and \(\max_i\|w_i\|^2\le Cm\).
The second-order normalization term consequently costs at most
\(Cm\Tr D_*^2\) after summing with weights \(m_i\).

The square terms in \eqref{eq:clearing-square-decomposition} are small
as well.  Pointwise Cauchy--Schwarz gives
\begin{equation}\label{eq:clearing-raw-squares}
 \sum_i|d_i^0|^2
 \le Cm\sum_i m_i|v_i|^2+Cm^2H_*
 \le Cm\sqrt{r_*mH_*}+Cm^2H_*.
\end{equation}
Write \(R_D=(\Id+D_*)^{-1}-\Id\).  Since \(R_D\) is Hermitian,
\[
 |w_iR_Dw_i^*|^2\le\|w_i\|^2w_iR_D^2w_i^*.
\]
Summing this inequality, and using
\(\|R_D\|_{\mathrm{HS}}^2\le C\Tr D_*^2\), gives
\begin{equation}\label{eq:clearing-normalization-squares}
 \sum_i|d_i-d_i^0|^2\le Cm\Tr D_*^2.
\end{equation}

All the errors just obtained are \(o(m^2)\).  More explicitly, after
division by \(m^2\), their bounds are constant multiples of
\[
 t^{\beta_w-e_0/2},\qquad
 t^{(\delta-e_0-\nu)/2},\qquad
 \varepsilon_t,\qquad t^{e_0/2},\qquad t^{e_0},\qquad H_*.
\]
Every exponent is positive under
\eqref{eq:clearing-drift-exponents}. This proves the claimed drift
of \(\mathcal P\).
\end{proof}

\subsection{Lowering the ridge and counting heats}
\label{subsec:clearing-ridge-count}

The procedure combines two motions: a fixed-ridge heat redistributes
weight while reducing $\mathcal P$ in conditional expectation, and a
ridge decrease raises the regularized weights toward the projection.
To control both with one quantity, define
\[
 \Phi=\mathcal P-10l_*\mathcal A.
\]
During a ridge decrease for which every $m_i\le2l_*$, each $dm_i$ is
nonnegative. Hence
\begin{equation}\label{eq:clearing-ridge-monotonicity}
 d\mathcal P=2\sum_i m_i\,dm_i\le4l_*\,d\mathcal A,
 \qquad d\Phi\le-6l_*\,d\mathcal A\le0.
\end{equation}
The same conclusion holds at the zero-ridge limit by continuity. At a
heat satisfying Lemma~\ref{lem:leverage-drift}, with $m=m_*\ge2l_*$,
the correction to the squared-weight drift obeys
\[
 10l_*|\EE(\mathcal A'-\mathcal A)|=o(l_*m)=o(m^2).
\]
Combining this with the drift of $\mathcal P$ gives, for small $t$,
\begin{equation}\label{eq:clearing-phi-drift}
 \EE(\Phi'-\Phi)\le-\frac12m_*^2.
\end{equation}
Thus a ridge decrease contributes no positive drift, while an admissible
heat with $2l_*\le m_*\le4l_*$ decreases $\Phi$ in conditional expectation by
at least $2l_*^2$.

Start with $\tau=l_*^{-2}$. Since $\mathcal C$ is a block of a unitary
matrix, $m_*\le\tau^{-1}=l_*^2<2l_*$. Whenever $m_*<2l_*$, decrease
$\tau$ until one of two events occurs. If the maximum diagonal of
$\Pi_{\Ran\mathcal C}$ is at most $2l_*$, declare completion at the
zero-ridge limit. Otherwise there is a unique positive ridge at which
$m_*=2l_*$, and the decrease stops there. For a nonzero row the displayed
derivative of $m_i(\tau)$ is strictly negative, which proves uniqueness.
The ridge is measurable: its comparison with any positive rational is
determined by the measurable value of $\max_i m_i(\tau)$ there. The
zero-ridge projection is measurable as the pointwise limit of the
positive-ridge matrices. At a heat, choose the least site among those
of maximum weight. This tie rule makes the pivot predictable.

\begin{lemma}[Number of heats]\label{lem:clearing-heat-count}
Consider an adapted procedure using the ridge rule above, with heats only at
positive ridge and at prestates satisfying Lemma~\ref{lem:leverage-drift}.
Stop the procedure at completion, at any prescribed failure, or at
a prescribed finite heat limit.  In particular, no further heat
is performed after one of the prestate hypotheses fails. The prescribed
stop is a stopping time for the revealed history. Let \(\sigma\) be any
almost surely finite stopping time at which this same procedure has reached
a prefix before that stop, with \(m_*\le4l_*\). Conditional on the retained
history at \(\sigma\), the expected number
\(H\) of further heats before this stop is at most
\[
 \EE H\le C\frac{r_*}{l_*}.
\]
There are constants \(c,C>0\), independent of the prefix, such that
\[
 \PP(H\ge n)\le C\exp\left(-c\,\frac{nl_*}{r_*}\right)
 \qquad(n\ge1).
\]
In particular, if the prescribed limit is
\(n_h=\lfloor t^{-1-e_h}\rfloor\) and \(e_h>e_0+\nu\), then
\begin{equation}\label{eq:clearing-heat-tail}
 \PP(H\ge n_h)
 \le C\exp\{-c\,t^{-(e_h-e_0-\nu)}\}.
\end{equation}
\end{lemma}

\begin{proof}
At this prefix, \(\mathcal P\le4l_*\mathcal A\), so \(\Phi\le0\).
At every later endpoint, including one reached by a failing heat,
\[
 \Phi\ge-10l_*\mathcal A\ge-10l_*r_*.
\]
The latter bound needs neither a leverage cap nor a diagonal cap:
it follows from the definition of ridge leverage and
\(\mathcal A\le|B|\).  The same statement holds at the projection
limit.  Sum the conditional drift \eqref{eq:clearing-phi-drift}
over the heats, including the last heat, and include the nonpositive
ridge changes from \eqref{eq:clearing-ridge-monotonicity}.
Stopping first at \(H\wedge k\) gives
\[
 2l_*^2\,\EE(H\wedge k)\le10l_*r_*.
\]
Letting \(k\) increase proves the expected-count estimate.

This estimate holds conditionally after every later reached prefix before stopping, with
the same constant.  Take a block length
\(k_{\mathrm{block}}=\lceil 2Cr_*/l_*\rceil\), enlarging \(C\) if necessary.
Conditional Markov inequality bounds the probability of a further
block of \(k_{\mathrm{block}}\) heats by \(1/2\).  On a prefix at which a stop has
already occurred the probability is zero.  Iterating over blocks
therefore gives
\(\PP(H\ge k k_{\mathrm{block}})\le2^{-(k-1)}\), which implies the stated exponential
bound after adjusting its constants.  Finally,
\[
 \frac{n_hl_*}{r_*}
 \ge c\,t^{-(e_h-e_0-\nu)}
\]
for sufficiently small \(t\), proving \eqref{eq:clearing-heat-tail}.
\end{proof}

For the chosen parameters,
\[
 \beta_w-\frac{e_0}{2}=.14\delta,\qquad
 \delta-e_0-\nu=(.56-10^{-5}/d)\delta,\qquad
 e_h-e_0-\nu=(.02-10^{-5}/d)\delta
\]
are all positive.  The leverage argument therefore bounds the
probability of exhausting the heat limit.  It remains to control the
other stopping conditions, so that the procedure reaches completion.

\subsection{The stopped trial and its local inputs}

We now state the local input that turns the range estimate into a trial
whose probability bound is independent of the torus volume. There is a
scale $\ell\asymp t^{-4}$ such that, as long as earlier heat pivots lie
within ordinary distance $\ell/2$ of $B$, the current collapsed distance,
and the distance after one further such deletion, satisfy
\begin{equation}\label{eq:clearing-distance-input}
 d_T(i,k)\ge\tfrac12\dist(i,k)-t^{-1.03}.
\end{equation}
The core lies in an ordinary-distance ball of radius $t^{-C_{\mathrm{loc}}}$,
for a fixed $C_{\mathrm{loc}}$. Every exterior site within distance
$\ell/2$ of $B$ which has not already been deleted is coupled and has a still
unexposed true residual. A predictable choice of such a site therefore
satisfies \eqref{eq:clearing-residual}. These are hypotheses about geometry
and recorded information; they contain no provisional bound on $K$ or
comparison for $D$.

Here is the selection consequence of the geometric input. By the
singular-value formula, a regularized weight is at most its projection
diagonal. For $\dist(i,B)\ge\ell/2$, Lemma~\ref{lem:clearing-range-localization}
and \eqref{eq:clearing-distance-input} bound that diagonal by
$C\exp(-ct^{-3})<2l_*$. Consequently
\begin{equation}\label{eq:clearing-proposed-location}
 m_x=m_*\ge2l_*\quad\Longrightarrow\quad \dist(x,B)<\ell/2.
\end{equation}
Starting with no earlier pivots, this implication proves location
inductively whenever the stated metric input is available under the
preceding-location hypothesis. It applies to the pivot chosen before a
first breaking heat. The input for the distance after that deletion will
also control its prospective output. Section~\ref{sec:assembly} supplies
this geometric input for the scheduled components.

For bookkeeping, keep a matrix $\hat h$ indexed by the sites active at trial
start.  A pair of entries is frozen immediately before the first of its two
endpoints is removed.  Only survivor--survivor entries receive a heat increment.
Put $\hat g=\hat h/\sqrt t$, and write $h^0,g^0$ for the trial-start matrices.
Thus each original row of $\hat h-h^0$, and the whole matrix viewed as a vector,
has a fixed index set throughout the trial.  The three distortion budgets are
\begin{equation}\label{eq:clearing-budgets}
\begin{split}
 \max_{i,k}|\hat g_{ik}-g^0_{ik}|&\le t^{.42\delta},\\
 \max_i\sum_k|\hat h_{ik}-h^0_{ik}|^2&\le t^{.25\delta},\\
 \sum_{i,k}|\hat h_{ik}-h^0_{ik}|^2&\le n_h t^{-.01\delta}.
\end{split}
\end{equation}

If $B$ is empty, declare DONE without changing the state. Otherwise the
trial starts at ridge $\tau=l_*^{-2}$. If the maximum ridge leverage is
below $2l_*$, lower the ridge to the first value where the maximum equals
$2l_*$; if the zero-ridge projection already has maximum leverage at most
$2l_*$, stop with outcome DONE. Otherwise choose an exterior site of maximum
regularized weight,
using a fixed ordering to break ties, perform one heat step, and recompute
the cross block and its leverages at the same scalar ridge.  Stop with outcome
FAIL if a distortion budget is exceeded, if a leverage exceeds $4l_*$, or if
$n_h$ heats have been performed.  Repeat the ridge lowering and heating rule
until one of these outcomes occurs. True values exposed during a trial
are part of the recorded history even when the state is later rolled back.

Call the state just before a heat performed by this rule a \emph{running
heat prestate}. Its three distortion budgets and leverage cap still hold,
fewer than $n_h$ heats have been performed, and the rule has selected a
positive ridge with $2l_*\le m_*\le4l_*$. These are exactly the conditions
under which the stopped rule performs another heat. The budget argument
below derives \eqref{eq:clearing-prestate} at these states from the starting
slack.
Write $\mathcal F_0$ for the retained history at trial start.

\begin{proposition}[Preparation of a small component]\label{prop:clearing-trial}
Assume \eqref{eq:clearing-broad}, the geometry and conditional residual
inputs just stated, and $1\le |B|\le r_*$.  At trial start suppose
\begin{equation}\label{eq:clearing-start}
 \max_i\sum_k|h^0_{ik}|^4\le .8H_*,\qquad
 \max_i|g^0_{ii}-r_i|=o(t^{\beta_w}),
\end{equation}
uniformly over the input states under consideration.  Assume also that every
true terminal principal matrix at the fixed energy is invertible.
For all sufficiently small $t$, the stopped trial is measurable and satisfies
\[
 \PP(\mathrm{FAIL}\mid\mathcal F_0)\le\exp(-t^{-c_2})
\]
for a fixed $c_2>0$.  The constants are independent of the torus volume and of
the component.  On DONE, fewer than $n_h$ exterior sites carry their true values and
are provisionally removed, the budgets \eqref{eq:clearing-budgets} hold, and
\begin{equation}\label{eq:clearing-projection-output}
 \max_{i\in R}(\Pi_{\Ran\mathcal C})_{ii}\le2l_*.
\end{equation}
Here $\mathcal C=S_{RB}$ is formed from the final active inverse before
completion of $B$. Thus DONE is the successful preparation outcome. The
subsequent deterministic core completion and the single successful commit
are specified with the clearing schedule in Section~\ref{sec:assembly}.
\end{proposition}

The uniform little-oh in \eqref{eq:clearing-start} is supplied by committed-state
assembly.  In the proof of Proposition~\ref{prop:clearing-trial} it is used only to obtain
\eqref{eq:clearing-prestate} from the entry budget.  We next verify every other
stopping probability.  All estimates are stopped at the first failure and
continued afterwards by zero increments; the step which first causes failure
is included.

\subsection{Distinct pivots and the distortion budgets}

The deterministic consequences of the budgets are useful before applying
concentration.  First,
\begin{equation}\label{eq:clearing-fourth-distortion}
 \sum_k|\hat h_{ik}-h^0_{ik}|^4
 \le\max_k|\hat h_{ik}-h^0_{ik}|^2
       \sum_k|\hat h_{ik}-h^0_{ik}|^2
 \le t^{1+1.09\delta}.
\end{equation}
The triangle inequality in $\ell^4$, the initial factor $.8$, and the positive
margin $.09\delta$ imply the fourth-power part of
\eqref{eq:clearing-prestate}.  The entry budget and
\eqref{eq:clearing-start} imply its lag part. Thus those provisional
hypotheses hold at every running heat prestate.

For a fixed original site $i$, sum over pivots $x$ for which $i$ is still active
immediately before the corresponding heat.  Each summand is evaluated at that
prestate.  Distinctness of the pivots and the previous entry budgets imply,
through a prospective breaking step,
\begin{equation}\label{eq:clearing-pivot-fourth}
 \sum_x F_{ix}^2
 \le C\sum_x|h^0_{ix}|^4
       +Cn_h t^{2+1.68\delta}
 \le C(H_*+t^{1+1.22\delta})\le CH_*.
\end{equation}
The first sum uses only the corresponding distinct columns of the initial row.
This is why a bound on the current row alone would not have been enough.
Cauchy--Schwarz now gives the pathwise predictable budget
\begin{equation}\label{eq:clearing-pivot-mass}
 B_i:=\sum_x F_{ix}\le C\sqrt{n_hH_*}
       \le Ct^{.27\delta}.
\end{equation}

For reference, the heat increment on surviving pairs is
\begin{equation}\label{eq:clearing-entry-increment}
 \Delta\hat g_{ik}=\frac{c_x}{\sqrt t}h_{ix}h_{xk},
 \qquad \Delta\hat h_{ik}=c_x h_{ix}h_{xk}.
\end{equation}
Entries with a removed endpoint have zero increment.  The conditional mean of
$c_x$ is $O(\sqrt t\,t^{\beta_w})$, its absolute value is bounded, and the
Ward row bound \eqref{eq:clearing-row-bound} is available at the prestate.
For a fixed entry, summing its conditional means in absolute value costs
\[
 Ct^{\beta_w}\sum_x\sqrt{F_{ix}F_{kx}}
 \le Ct^{\beta_w}\sqrt{B_iB_k}.
\]
The conditional variance sum is at most
\[
 \frac Ct\sum_xF_{ix}F_{kx}\le CH_*/t,
\]
by \eqref{eq:clearing-pivot-fourth}, and each jump is at most
$C\sqrt{H_*/t}=Ct^{\delta/2}$.

For a fixed row regarded as a vector, its increment has norm at most
$C|h_{ix}|$; the same estimate with $C$ replaced by
$C\sqrt t\,t^{\beta_w}$ bounds its conditional mean.  Consequently the total
mean norm, predictable variance, and maximum jump are bounded respectively by
\[
 C\sqrt t\,t^{\beta_w}\sqrt{n_hB_i},\qquad
 CB_i,\qquad CH_*^{1/4}.
\]
For the whole matrix regarded as a Frobenius vector, the increment norm is at
most $C\|h_{\cdot x}\|_2^2\le C$, where the dot denotes the full current
active column. Its corresponding three bounds are
\[
 C\sqrt t\,t^{\beta_w}n_h,\qquad Cn_h,\qquad C.
\]
The conditional means here are computed before the query: the endpoint mask
is predictable, so the scalar $\EE c_x$ can indeed be factored out of the
increment.  This mask is a contraction in the row and Frobenius norms and
does not change any bound.  Center each increment at its conditional mean.
For a complex vector $z$, its Hermitian dilation
\[
 \mathscr D(z)=\begin{pmatrix}0&z^*\\z&0\end{pmatrix}
 \quad\hbox{satisfies}\quad
 \mathscr D(z)^2=\diag(\|z\|^2,zz^*)\le\|z\|^2\Id.
\]
Centering cannot increase the conditional second moment in matrix order.
Thus the displayed squared-norm bounds give predictable \emph{scalar}
variance envelopes in the fixed original coordinate spaces. For entries and
rows their pathwise sums follow from \eqref{eq:clearing-pivot-fourth} and
\eqref{eq:clearing-pivot-mass}; for the Frobenius vector the Ward bound and
the limit of $n_h$ heats give the sum directly. All include the candidate
first breaking heat.
This verifies the precise hypotheses of
Lemma~\ref{lem:matrix-concentration}, rather than requiring a matrix-valued
variance theorem.

Here are the margins at the required thresholds.  For entries, use
$a=t^{.42\delta}$; the total drift is $O(t^{.63\delta})$ and
\[
 \frac{a^2}{V}\ge ct^{-.16\delta},\qquad
 \frac aJ\ge ct^{-.08\delta}.
\]
For row norms, use $a=t^{.125\delta}$; the drift is
$O(t^{.265\delta})$ and $a^2/V\ge ct^{-.02\delta}$, with a larger
margin for $a/J$.  For the Frobenius norm, use
$a=\sqrt{n_h}\,t^{-.005\delta}$; its drift divided by $a$ is
$O(t^{.135\delta})$, and $a^2/V\ge ct^{-.01\delta}$.
Thus all three budget crossings have stretched-exponentially small
probability after restriction to a polynomial number of coordinates.  We
justify that restriction below, before drawing the conclusion.

\subsection{A two-sided bound on leverage changes}

The drift calculation controls the number of heats, but it does not by itself
keep individual leverages below $4l_*$.  We prove that cap separately.  For each
original exterior site $i$, record the fixed-ridge heat change
$d_i=m'_i-m_i$ whenever $i$ survives the heat.  Set the recorded change to zero
at its own deletion and after deletion.  Ridge adjustments are not included.
More precisely, number the heats by $k$. While the trial is running, let
$\mathcal F_{k-1}$ be the full retained history immediately before the
residual query for heat $k$, including the already selected pivot and ridge
$\tau_k$. Put
\[
 Z_{k,i}={\bf1}_{\{i\text{ survives heat }k\}}
       \{m_i^{\mathrm{post}}(\tau_k)-m_i^{\mathrm{pre}}(\tau_k)\}.
\]
The survival indicator is predictable, because the pivot is chosen before
its residual is queried. After stopping, the formal remaining increments
are zero and $\mathcal F_k$ stays equal to the history at stopping. The
actual retained filtration continues through rollback, marking, and later
trials.
Thus $Z_{k,i}-\EE(Z_{k,i}\mid\mathcal F_{k-1})$ is one adapted martingale
stream even though $\tau_k$ can vary with $k$.  Every estimate below is uniform
in that predictable positive ridge; no discretization or union over its
possible values is needed.  The notation $d_i$ denotes $Z_{k,i}$ at the heat
under consideration.

Write $\varepsilon_t=t^{\beta_w}+\sqrt t+r_*\sqrt{H_*}$.  The same expansion
of the Gram normalization used in Lemma~\ref{lem:leverage-drift} gives, at a
running heat prestate,
\begin{equation}\label{eq:clearing-one-row-bias}
 |\EE(d_i\mid\mathcal F)|
 \le Cl_*\left[
 \sqrt t\,t^{\beta_w}|v_i|+|v_i|^2
 +l_*\varepsilon_t+\frac{l_*|v_i|}{\sqrt t}
 +\frac{l_*^2}{t}\right],
\end{equation}
\begin{equation}\label{eq:clearing-one-row-jump}
 |d_i|\le Cl_*\left(|v_i|+\frac{l_*}{\sqrt t}\right).
\end{equation}
For completeness, the first two terms of
\eqref{eq:clearing-one-row-bias} are the conditional mean of the raw row-norm
change.  In the first-order normalization, keeping the old row gives
$Cl_*\|\EE D_*\|_1\le Cl_*^2\varepsilon_t$.  Replacing it by
the new row costs $Cl_*|v_i|\|D_*\|\le
Cl_*^2|v_i|/\sqrt t$.  The quadratic remainder is at most
$Cl_*\operatorname{Tr}D_*^2\le Cl_*^3/t$.
The pathwise estimate follows by using $|c_x|\le C$ and
$\|D_*\|\le Cl_*/\sqrt t$ instead of their means.

Since the weights $p_i$ are bounded above and below,
\[
 \sum|v_i|^2\le CB_i,\qquad
 \sum|v_i|\le C\sqrt{n_hB_i}.
\]
The sum of the absolute conditional biases in
\eqref{eq:clearing-one-row-bias}, divided by $l_*$, is therefore bounded by
\begin{equation}\label{eq:clearing-bias-margins}
 C\left[
 t^{.265\delta}+t^{.27\delta}+t^{.34\delta}
 +t^{1/2-.02\delta}+t^{1/2+.48\delta-\nu}
 +t^{.345\delta}+t^{.42\delta}\right]=o(1).
\end{equation}
The conditional variance sum is at most
\begin{equation}\label{eq:clearing-leverage-variance}
 Cl_*^2\left(B_i+n_hl_*^2/t\right)
 \le Cl_*^2t^{.27\delta}.
\end{equation}
The jump divided by $l_*$ is at most
$C(H_*^{1/4}+l_*/\sqrt t)$, a positive power of $t$.

Apply both signs of the maximal martingale inequality at a sufficiently small
fixed multiple of $l_*$.  After adding the absolute bias bound
\eqref{eq:clearing-bias-margins}, except with stretched-exponentially small
probability we have
\begin{equation}\label{eq:clearing-two-sided}
 \sup_s\left|\sum_{k\le s}d_i^{(k)}\right|\le l_*/2
\end{equation}
at every prefix and every monitored row.  The absolute value in
\eqref{eq:clearing-two-sided} is essential: it bounds the change over any
interval of heat steps by $l_*$.

Suppose a surviving row first exceeded $4l_*$.  Immediately after the most
recent ridge lowering, or at initialization if there was none, its leverage
was at most $2l_*$.  Let $k_0$ be the number of heats preceding that adjustment and $s$ the
current number.  The ridge has been held constant for heats $k_0+1,\ldots,s$,
and the row has survived all of them.  Its actual change is therefore exactly
\[
 \sum_{k=k_0+1}^s Z_{k,i}
   =\sum_{k\le s}Z_{k,i}-\sum_{k\le k_0}Z_{k,i}.
\]
The two prefix sums use the single varying-ridge stream; their difference
is the desired constant-ridge interval sum.  Equation~\eqref{eq:clearing-two-sided}
bounds it in absolute value by $l_*$, giving leverage at most $3l_*$,
a contradiction. This argument includes a prospective cap-breaking heat and
uses only running heat prestates in its concentration estimates.

\subsection{Finite monitoring and conclusion}

It remains to make the stopped probability estimates independent of the
ambient torus size. The location implication
\eqref{eq:clearing-proposed-location} keeps every selected pivot within
$\ell/2$ of $B$, including the pivot of a prospective first failure.
The residual hypothesis therefore applies at every heat.

Choose a fixed ordinary-distance cube $Q$ containing $B$, its
$\ell$-neighborhood, and a further buffer of width $t^{-5}$.  By the bound
on the diameter of $B$, the number of sites in $Q$ is polynomial in $1/t$.
Outside $Q$, projection decay rules out a leverage-cap crossing
deterministically, even at a prospective failing output.  For the distortion
processes, the collapsed-distance resolvent estimate gives
\[
 |h_{ik}|\le Ct^{-1/2}e^{-ctd_T(i,k)}.
\]
Each heat increment with an endpoint outside $Q$ has an exponentially small
factor connecting that endpoint to the pivot.  The other factor has bounded
row square sum by \eqref{eq:clearing-row-bound}.  More explicitly, the
squared Frobenius norm of the portion of one increment with at least one
endpoint outside $Q$ is at most
\[
 C\left(\sum_{i\notin Q}|h_{ix}|^2\right)
   \left(\sum_k|h_{xk}|^2\right).
\]
The first sum has a dimension-dependent bound which is still independent
of the torus volume.  An ordinary shell of integer radius $q$ has at most
$C_d(1+q)^{d-1}$ sites.  For an endpoint outside $Q$ its distance from the
pivot is at least $t^{-5}$, so \eqref{eq:clearing-distance-input} gives
$d_T(i,x)\ge\dist(i,x)/4$ for small $t$.  Therefore
\[
 \sum_{i\notin Q}|h_{ix}|^2
 \le C_dt^{-1}\sum_{q\ge t^{-5}}(1+q)^{d-1}e^{-ctq}
 \le C_dt^{-d-1}e^{-c't^{-4}}.
\]
In the last step, extract $e^{-ct^{-4}/2}$ and bound the remaining
exponential moment by $C_dt^{-d}$.  Summing at most $n_h$ increments preserves an
exponential bound by the triangle inequality.  The same reasoning gives the
entry and individual-row tails.  It uses no factor proportional to the torus
volume.

We may therefore apply the martingale inequalities to entries and row vectors
compressed to $Q$, and to the matrix compressed to $Q\times Q$, and then add
the deterministic tails.  Compressions do not increase the conditional
square bounds.  If the core lies in a ball of radius $t^{-C_{\mathrm{loc}}}$, then
\[
 |Q|\le C_d(t^{-C_{\mathrm{loc}}}+t^{-4}+t^{-5})^d.
\]
Thus the largest vector-dilation dimension is at most $1+|Q|^2$, and there
are at most $C|Q|^2$ entry, row, and leverage tests.  These are fixed powers
of $1/t$ for each fixed $d$.  Take the martingale thresholds to be one quarter of the
budget thresholds.  Their predictable biases and deterministic tails are
smaller than another quarter for sufficiently small $t$.  The strict margins
proved above imply that the probability of a distortion crossing or leverage
crossing is at most $\exp(-t^{-c})$ for some fixed $c>0$.

Finally, Lemma~\ref{lem:clearing-heat-count} bounds the probability of reaching
$n_h$ heats before another stop by
$C\exp[-c t^{-(e_h-e_0-\nu)}]$.  Because
$e_h-e_0-\nu=(.02-10^{-5}/d)\delta>0$, this is of the same required form after
decreasing the exponent.  There are no other failure modes.  The procedure
uses measurable spectral functions, a measurable threshold ridge and a fixed
tie rule; each real update is measurable as well.  It must stop after at most
$n_h$ heats.  On DONE the zero-ridge projection satisfies
\eqref{eq:clearing-projection-output} by the stopping rule, and the budgets
hold.  This proves Proposition~\ref{prop:clearing-trial}.

On DONE the preparation has supplied bounded accumulated heat distortion and
a small diagonal for the final full-range projection. Lemma~\ref{lem:core-completion}
turns that projection bound into deterministic control of the true core
completion, and Corollary~\ref{cor:core-decay} gives its spatial decay.
The next section specifies the successful commit and the rollback of a
failed preparation, and accumulates successful operations using comparisons
only at committed states.
